% Separate post-CP17 candidate. No unrestricted tame least-index conclusion.
\subsection{The remaining tame ambiguity cannot have two CM formal models}

Here a potentially supersingular height-two group is called
\emph{potentially CM} if, over some finite extension on which it
has good formal reduction, its rational endomorphism algebra
contains a quadratic field over $\mathbf Q_p$. This requires
endomorphisms defined over a finite extension; it is not the
endomorphism algebra of the constant etale group over an
algebraic closure of the generic field.

\begin{lemma}[Quadratic endomorphisms and a nonsquare determinant]
\label{lem:B-RD52-CM}
\status{Proved; complete root and obstruction/crystal proof, original-source and exact-code audits passed.}
Let $p\geq5$, and let $K$ be complete discretely valued of mixed
characteristic $(0,p)$, with algebraically closed residue field,
$e_K=p-1$, and $\mu_p\subset K$. Let $M$ be a potentially CM,
potentially supersingular height-two group. Write its determinant
as $\chi_p\psi$, where $\psi$ is its finite special determinant.
If $\psi$ is nonsquare as a $\mu_{p-1}$-valued character, then
its CM field is the unramified quadratic field $E=\mathbf Q_{p^2}$,
all its finite-extension endomorphisms are defined over $K$, and its Tate lattice
is free of rank one over $\mathcal O_E$.
\end{lemma}
\begin{proof}
Pass to a faithful tame good field $L/K$, as in
Lemma~\ref{lem:B-RD49-inertia}. Tate's full faithfulness
\cite[Theorem~4 and Corollaries~1--2, p.~181]{Tate1967PDivisible}
extends a generic finite-extension endomorphism to the good
Barsotti--Tate model; for a quasi-endomorphism, first clear a
power of $p$. Reduction of homomorphisms between finite-height
one-dimensional formal groups is injective, as in
\cite[Proposition~4.2 and its full proof]{HopkinsGross1994Bundles}.
Indeed, if a nonzero homomorphism had all coefficients divisible
by $\pi$, divide by their minimum positive valuation, say $\pi^r$.
Reducing the identity $f\circ[p]=[p]\circ f$ modulo
$\pi^{r+1}$ would give a nonzero series composed with the
nonzero special $[p]$ on the left and zero on the right.
This is impossible. The assertion also holds rationally.

Any endomorphism defined after a further finite extension of
$L$ therefore descends to $L$: after a common finite Galois
extension, its Galois conjugates have the same reduction, since
the residue field is algebraically closed and the model over
$L$ has good reduction. Injectivity makes the conjugates equal.
The rational endomorphism algebra over $L$ embeds both in the
height-two division algebra $D$ of the special fibre and in
$M_2(\mathbf Q_p)$ on the Tate module. It is a division algebra,
and the two-dimensional Tate module is a vector space over it;
its dimension over $\mathbf Q_p$ is consequently at most two.
Under the CM hypothesis it is a quadratic field $E$.

The cyclic finite special inertia group $\Gamma\subset\mathcal O_D^\times$
normalizes $E$. If $E$ is ramified, its prime-to-$p$ roots of
unity have order dividing $p-1$, and
$N_D(E)/E^\times$ has order two. Hence $|\Gamma|\mid2(p-1)$.
If $p\equiv3\pmod4$, every such root in
$\mathbf F_{p^2}^\times$ has square reduced norm: its
$(p^2-1)/2$ power is one. This makes $\psi$ square.
If $p\equiv1\pmod4$, use the generic normalizer instead.
On the one-dimensional $E$-space $V_p(M)$, its elements are
$a$ or $a\iota$, where $\iota$ is the nontrivial field
automorphism; their determinants are respectively
$N_{E/\mathbf Q_p}(a)$ and $-N_{E/\mathbf Q_p}(a)$.
A Galois element preserves the Tate lattice, so this determinant
is a unit and $a\in\mathcal O_E^\times$. The residue of its
norm is a square for a ramified quadratic field, and $-1$ is
also a square when $p\equiv1\pmod4$. Since
$\chi_p\equiv1\pmod p$, the residue of $\psi$ is again
always square. Thus ramified $E$ is impossible.

For unramified $E$, write
$D=E\oplus E\Pi$, with $\Pi^2=p$ and
$\Pi a=\sigma(a)\Pi$, where $\sigma$ is the nontrivial
automorphism of $E/\mathbf Q_p$. The noncentral coset of $N_D(E)$ consists
of $E^\times\Pi$, whose elements have odd $D$-valuation.
It has no units. Therefore $\Gamma$ centralizes $E$,
so the whole generic Galois action centralizes $E$.

Choose a Galois element $g$ for which $\psi(g)$ is nonsquare.
Then $\rho(g)\in E^\times$ has unit norm with nonsquare
residue, and hence its residue is not in $\mathbf F_p$.
It follows that $\mathbf Z_p[\rho(g)]=\mathcal O_E$:
its quadratic characteristic polynomial has irreducible,
separable reduction. The Tate lattice, stable under $\rho(g)$,
is thus stable under $\mathcal O_E$, and is free of rank one
over this discrete valuation ring.
\end{proof}

\begin{lemma}[The integral CM model and its full first wild layer]
\label{lem:B-RD52-LT}
\status{Proved; complete root and obstruction/crystal proof, original-source and exact-code audits passed.}
In the nonsquare case of Lemma~\ref{lem:B-RD52-CM}, choose the
orientation of $E$ on the tangent line. Let $F_0$ be the
Lubin--Tate formal $\mathcal O_E$-module with
$[p](T)=pT+T^{p^2}$, and let $\theta$ denote its Tate
character after base change to $K$. Then
\[
                  T_p(M)\simeq\mathcal O_E(\theta\epsilon)
                                                        \tag{B-RD52.1}
\]
as a $\mathbf Z_p$-representation, for a finite tame character
$\epsilon:G_K\to\mu_{p^2-1}$. Over every common faithful
tame good field, the image of wild inertia under $\theta$
modulo $p^2$ contains all of
$(1+p\mathcal O_E)/(1+p^2\mathcal O_E)$.
\end{lemma}
\begin{proof}
Tate's full faithfulness extends the integral
$\mathcal O_E$-action on the Tate lattice to the good formal
model; see \cite[Theorem~4 and Corollaries~1--2, p.~181]
{Tate1967PDivisible}. We include the uniqueness argument for
this model over an arbitrary algebraically closed perfect
residue field, rather than silently identifying that field
with $\overline{\mathbf F}_p$.

Here the integral conjugacy can be checked without passing
through a rational conjugator. Put $N=p^2-1$ and let $v$
be the image of a primitive $N$th root under the special
CM action. After possibly interchanging the two embeddings
of $E$, its residue in $\mathcal O_D/(\Pi)$ equals that of
the corresponding standard root $\lambda$. Then
\[
 x=\sum_{i=0}^{N-1}v^i\lambda^{-i},\qquad
 vx=x\lambda,\qquad x\equiv N\pmod\Pi.
\]
Thus $x$ is an integral unit and conjugates the whole
$\mathcal O_E$-action to the standard action on the
reduction of $F_0$.

The height-two universal deformation ring is $W(k)[[u]]$;
the arbitrary-residue-field assertion and the degree-$p$
coefficient are \cite[Proposition~5.7 and equations~(5.3)--(5.8);
Proposition~12.10 and its full proof]{HopkinsGross1994Bundles}.
Let $\lambda$ be a primitive $(p^2-1)$st root. Its special
automorphism induces an action on this ring whose tangent
eigenvalue is $a=\lambda^{p-1}\ne1$; this follows directly
by conjugating the degree-$p$ term by $T\mapsto\lambda T$.
Choose the parameter to vanish at the canonical deformation
$F_0$. Averaging it against this finite prime-to-$p$ action
gives a parameter $u_*$ on which the action is exactly
$u_*\mapsto a u_*$. An $\mathcal O_E$-linear deformation
is fixed by this action, so its classifying homomorphism
annihilates $u_*$, since $a-1$ is a unit. It is therefore
the base change of $F_0$. This proves the required integral
uniqueness. Changing the chosen orientation conjugates both
$\theta$ and $\epsilon$ by the unramified field automorphism,
and does not change the underlying rank-two special character.

The Tate character of the good model is consequently
$\theta|_{G_L}$. The quotient by $\theta$ over $K$ is an
$\mathcal O_E^\times$-valued character trivial on $G_L$,
and hence a finite tame character $\epsilon$. This gives
(B-RD52.1).

For the wild assertion put $q=p^2$ and $K_0=W(k)[1/p]$.
The elementary proof of
\cite[Theorem~2(b),(d),(e) and full proof, pp.~383--384]
{LubinTate1965FormalCM} applies after this residue extension.
Indeed, the primitive $p^2$-torsion polynomial is
\[
                 \Phi_2(T)=(pT+T^q)^{q-1}+p.           \tag{B-RD52.2}
\]
It remains Eisenstein over $W(k)$, of degree $q(q-1)$.
Every primitive root generates its entire division field:
the other torsion points are its $\mathcal O_E/p^2$ multiples.
The resulting Galois group is exactly
$(\mathcal O_E/p^2)^\times$. Its wild Sylow subgroup is
$(1+p\mathcal O_E)/(1+p^2\mathcal O_E)$.
Since $L/K_0$ is tame, its intersection with this division
field is fixed by that Sylow subgroup. The subgroup therefore
survives unchanged after base change to $L$.
\end{proof}

\begin{theorem}[Exact threshold in the potentially CM formal subclass]
\label{thm:B-RD52}
\status{Proved; complete root and obstruction/crystal proof, original-source and exact-code audits passed.}
Let $K$ be complete discretely valued of mixed characteristic
$(0,p)$ with perfect residue field, where $p\geq5$,
$e_K=p-1$, and $\mu_p\subset K$. Two
potentially CM, potentially supersingular height-two groups
with isomorphic $p^2$-torsion have equal requested special
characters. The level $p$ does not suffice within this same
subclass. Thus the least manuscript index in this subclass
is exactly one.

In particular any remaining pair in Theorem~\ref{thm:B-RD49}
with unequal characters and isomorphic $p^2$-torsion must
have at least one formal model without quadratic
finite-extension endomorphisms. This does not exclude
non-CM deformations in either of the two patterns of
Corollary~\ref{cor:B-RD51-p5}.
\end{theorem}
\begin{proof}
Pass to the completed maximal unramified base. The determinant
comparison in Theorem~\ref{thm:B-RD49} makes the two finite
determinants equal. If that common character is square,
Proposition~\ref{prop:B-RD50-square} already proves equality.
In the nonsquare case, the preceding lemmas give integral
identifications $T_p(M_i)=\mathcal O_E(\theta\epsilon_i)$,
with the same oriented $\theta$ after changing an
$\mathcal O_E$-orientation if needed.

Let $A$ be an intertwiner modulo $p^2$. On wild inertia in
a common tame good field, both representations contain
every multiplication operator $1+pb$, $b\in\mathcal O_E$.
Hence the reduction $\bar A$ commutes with all multiplication
operators from $\mathbf F_{p^2}$, and is itself multiplication
by a nonzero element of that field. Reducing the intertwining
identity for any $g\in G_K$ modulo $p$ now gives
$\bar\epsilon_1(g)=\bar\epsilon_2(g)$, since $\bar A$
commutes with all field scalars and $\bar\theta(g)\ne0$.
Reduction is injective on these prime-to-$p$ roots of unity;
therefore $\epsilon_1=\epsilon_2$. The special characters
are equal.

For the sharpness assertion, take $k=\overline{\mathbf F}_p$,
$K_0=W(k)[1/p]$, and $K=K_0(\mu_p)$. Choose a uniformizer
$\pi\in K$ with $\pi^{p-1}=-p$: the cyclotomic extension
has tame degree $p-1$, and over the strictly henselian
$K_0$ the unique tame extension of that degree is
$K_0((-p)^{1/(p-1)})$. Equivalently, remove the unit in a
uniformizer's $(p-1)$st power by Hensel's lemma.
Write $\omega$ for
the Teichmuller tame character given by a root
$\pi^{1/(p^2-1)}$. The first torsion equation for $F_0$ is
$T^{p^2-1}=-p$, so
$\bar\theta=\bar\omega^{p-1}$.
Twist $F_0$ generically by the two finite
$\mathcal O_E$-linear characters
\[
            \epsilon_1=\omega,\qquad
            \epsilon_2=\omega^{p^2-p+1}.               \tag{B-RD52.3}
\]
These twists are effective on the generic p-divisible groups
by finite Galois descent, and have good CM reduction after
a tame extension. Their mod-$p$ representations are,
respectively, multiplication by $\bar\omega^p$ and by
$\bar\omega$, which are isomorphic over $\mathbf F_p$
via field Frobenius.
Their finite special eigenvalue exponents are respectively
$\{1,p\}$ and $\{2-p,2p-1\}$ modulo $p^2-1$.
Here the exponents use the covariant convention in
Proposition~\ref{prop:B-RD4-dictionary}:
$D_{\mathrm{cov}}(G)[1/p]=D_{\cris}(V_pG)(-1)$.
The good Lubin--Tate factor has trivial finite special inertia,
the finite twist contributes $\epsilon_i$, and the period
Tate twist has trivial finite descent inertia. Thus these
are the actual displayed covariant exponents, not merely
an assertion up to inversion.
The second first exponent is neither $1$ nor $p$ modulo
$p^2-1$ for $p\ge5$, so the special characters differ.
Both finite determinants have the same primitive norm.
This proves failure at $p$ and the exact restricted index.
\end{proof}

\begin{remark}[Certificate and exact restriction]
\label{rem:B-RD52-scope}
\status{Proved under the stated CM restriction; no unrestricted least-index claim.}
The retained exact program
\path{code/reduction_descent/verify_tame_CM_gate.py}
checks the symbolic quadratic-field centralizer, all
$260{,}070$ two-by-two matrices in its six small-prime cases,
$75{,}984$ finite-field norm cases, $2{,}264$ ramified-norm
cases, $1{,}014$ Eisenstein coefficients, and the displayed
lower character pair for every prime from $5$ through $101$.
It also verifies $75{,}960$ terms in the unit-conjugacy telescoping identity.
Its final successful author runtime was $0.3032921849953709$ seconds.
Independent root and obstruction/crystal reruns took
$0.33256562399765244$ and $0.3663983289952739$ seconds,
respectively, and preserved the author output.
The first run used structural equality of unexpanded symbolic
matrix entries; that failed version and output are retained,
and the corrected comparison expands each difference exactly.
These computations verify finite algebra. The integral
deformation, endomorphism and division-field arguments are
the proofs above. They do not imply that every remaining
tame formal model is potentially CM, nor do they replace
the unresolved non-CM cases by elliptic curves.
\end{remark}
